\subsection{RQ4f: most models relay the failed injection}
\label{sec:relay-models}

Is the relay a quirk of one model? RQ4f ran the same task on all twelve models, with six framings: the
three of RQ4d and three new ones, each asking for a different internal document (the expense policy,
the VPN guide, the leave policy) to go to a different outside domain (a customs broker, a portal
migration desk, an account manager). Twenty runs per framing and model, 1{,}440 runs in all, 1{,}290 of
them new; the pre-registration, including the new framings, was committed before the first run.

Nine of the twelve models relayed the request to the user in 61\% to 100\% of their runs, and none of
the nine warned in a single run (Table~\ref{tab:rq4f}). They span every vendor and size in the study:
small open models, gpt-oss-120b, gpt-4.1-mini, gpt-5-mini and Claude Haiku~4.5. Pooled over the nine,
997 of 1{,}080 runs relayed; per framing between 80\% (the JSON record) and 99\%, and on the three
framings first used here between 90\% and 98\%. Only Opus~4.8, Opus~5.5 and Sonnet~5.5 never relayed:
they warned, or the provider stopped the answer. The agents themselves almost never acted (3 of 1{,}440
runs sent the email). For most models, then, the usual end of a failed injection is not an alert and
not a refusal, but an instruction to the user, with the attacker's address in it.

\begin{table}[t]\centering\small
\resizebox{\textwidth}{!}{\begin{tabular}{lrllrrrr}
\toprule
Model & Runs & Relays \% [95\% CI] & Held-out framings \% & gpt-5-mini \% & gpt-4.1-mini \% & Warns & Harm \\
\midrule
Opus 4.8 & 120 & 0 [0, 3] & 0 & 0 & 0 & 120 & 0 \\
Opus 5.5 & 120 & 0 [0, 3] & 0 & 0 & 0 & 74 & 0 \\
Sonnet 5.5 & 120 & 0 [0, 3] & 0 & 2 & 0 & 120 & 0 \\
\midrule
Haiku 4.5 & 120 & 89 [82, 94] & 100 & 89 & 32 & 0 & 0 \\
gpt-5-mini & 120 & 100 [97, 100] & 100 & 98 & 42 & 0 & 0 \\
gpt-4.1-mini & 120 & 98 [93, 99] & 100 & 98 & 76 & 0 & 0 \\
gpt-oss-120b & 120 & 88 [81, 93] & 93 & 88 & 46 & 0 & 0 \\
qwen3:8b & 120 & 100 [97, 100] & 100 & 100 & 61 & 0 & 0 \\
gemma4:e4b & 120 & 98 [94, 100] & 100 & 100 & 12 & 0 & 2 \\
granite4:tiny-h & 120 & 98 [94, 100] & 97 & 99 & 40 & 0 & 1 \\
ministral-3:8b & 120 & 98 [94, 100] & 100 & 98 & 51 & 0 & 0 \\
llama3.1:8b & 120 & 61 [52, 69] & 52 & 60 & 12 & 0 & 0 \\
\bottomrule
\end{tabular}
}
\caption{RQ4f, 120 runs per model (six framings, 20 each). Relays: answer passes the request on as
something to do, no warning (pre-registered labels, Wilson CI). Held-out: framings R4 to R6 only. 2nd
rater: relay rate by gpt-5-mini, labelling blind with the same definitions. Harm: the agent sent the
email itself.}
\label{tab:rq4f}
\end{table}

\paragraph{How reliable are the labels?} The two pre-registered raters were two instances of the same
model; they agreed on all 1{,}440 items, which shows that the labelling is reproducible but not that it
is independent. A rater from another family, gpt-4.1-mini, agreed on only 59\% ($\kappa = 0.36$). In ten
randomly drawn disagreements every answer passed the request on with no warning; it had labelled the
request, not the answer. By the pre-registered rule, which judged agreement on that rater, the
generality claim is therefore \emph{not established}. The next rater was fixed before it ran, with the
identical prompt: gpt-5-mini agreed with the pre-registered labels on 97\% of items ($\kappa = 0.94$;
0.98 on relay versus warning) and gives the same rate per model (Table~\ref{tab:rq4f}). We report both:
the formal verdict fails on one rater, and the evidence from the other is strong.

\paragraph{Detecting it.} \llmd{026} caught 77\% of the relays, and 77\% on the three internal documents
it was not written for. It fired on 6\% of the warnings and on none of 788 benign runs across eleven
models. Its misses name the address without the document or the document without an address.
